\section{Results}\label{sec:results}

\subsection{Cohorts}
After exclusions we analysed \NGradedTotal{} patients with a recorded grade in \NGradedCohorts{}
cohorts, plus \NTcga{} ungraded TCGA patients for description (Table~\ref{tab:cohorts};
participant flow in Supplementary Table~\ref{tab:flow}). The cohorts differed widely
(Figure~\ref{fig:grade}). Median age ranged from \AgeMbc{} years in MBCproject to \AgeMetabric{} years
in METABRIC. Grade~3 made up \GThreeMetabric\% of METABRIC, \GThreeMsk\% of MSK-2018, which is
enriched for advanced disease, and \GThreeDuke\% of Duke. Compared with every open cohort except the
triple-negative BCCRC series, the Dammam Saudi series~\citep{kussaibi2021saudi} had fewer ER-positive
(57.3\%) and PR-positive (45.0\%) tumours; only Sanger had more HER2-positive tumours (Dammam
25.3\%). Even so, the Dammam grade distribution (8.7/54.5/36.7\%) had more grade~2 and less grade~3
than METABRIC.

\begin{table}[htbp]\centering\small
\caption{Characteristics of the open cohorts after harmonisation and of two published Saudi hospital
series. Percentages are of patients with the variable recorded (missingness in
\texttt{results/e1\_cohorts.csv}). MSK = MSK-2018; MBCP = MBCproject; Broad = Mexico and Vietnam;
dev/ext/desc = development, external validation, descriptive. $^\dagger$Saudi nationals, King Fahd
University Hospital, Dammam, 1998--2017~\citep{kussaibi2021saudi}. $^\ddagger$King Fahad Medical
City, Riyadh, 2016--2019 (90.1\% Saudi nationals; 39.8\% stage IV)~\citep{durrani2025healthcare};
receptor percentages recomputed from the published counts.}\label{tab:cohorts}
\tabfit{\begin{tabular}{lrrrrrrrrrr}
\toprule
 & METABRIC & MSK & Duke & MBCP & Sanger & Broad & BCCRC & TCGA & Dammam$^\dagger$ & Riyadh$^\ddagger$ \\
Role & dev & ext & ext & ext & ext & ext & desc & desc & anchor & anchor \\
\midrule
Patients analysed, $n$ & 2386 & 1573 & 637 & 207 & 100 & 59 & 63 & 1072 & 286 & 283 \\
Age, median (IQR), y & 61 (51--70) & 52 (43--61) & 52 (45--61) & 43 (38--53) & 55 (44--66) & 52 (42--58) & 56 (48--63) & 58 (49--67) & 51 & -- \\
Age $\le$ 50 y, \% & 23.4 & 45.7 & 42.4 & 70.0 & 43.0 & 45.8 & 33.3 & 30.4 & -- & 45.2 \\
Grade 1, \% & 8.9 & 5.7 & 17.7 & 12.1 & 8.0 & 6.8 & 6.3 & -- & 8.7 & 4.2 \\
Grade 2, \% & 40.9 & 26.7 & 49.8 & 50.2 & 41.0 & 49.2 & 4.8 & -- & 54.5 & 47.0 \\
Grade 3, \% & 50.2 & 67.6 & 32.5 & 37.7 & 51.0 & 44.1 & 88.9 & -- & 36.7 & 48.8 \\
IDC / NST, \% & 76.7 & 81.8 & 89.0 & 83.0 & 79.2 & 96.6 & 92.1 & 77.2 & 91.3 & -- \\
ILC, \% & 7.9 & 11.7 & 9.9 & 10.2 & 6.2 & 0.0 & 0.0 & 20.2 & 7.0 & -- \\
Mixed ductal--lobular, \% & 11.2 & 4.9 & 0.0 & 5.8 & 9.4 & 1.7 & 1.6 & 0.0 & -- & -- \\
Other special type, \% & 4.2 & 1.5 & 1.1 & 1.0 & 5.2 & 1.7 & 6.3 & 2.6 & 1.7 & -- \\
ER positive, \% & 73.9 & 84.6 & 75.5 & 84.1 & 79.0 & 68.6 & 4.8 & -- & 57.3 & 62.6 \\
PR positive, \% & 52.1 & 72.2 & 64.8 & 74.7 & 60.0 & 58.8 & 1.6 & -- & 45.0 & 54.1 \\
HER2 positive, \% & 12.5 & 11.5 & 17.6 & 23.6 & 30.0 & 17.4 & 0.0 & -- & 25.3 & 36.1 \\
T1 ($\le$20 mm), \% & 43.8 & 49.2 & 44.9 & -- & -- & -- & 42.6 & 25.8 & -- & -- \\
T2 (21--50 mm), \% & 51.0 & 34.4 & 42.7 & -- & -- & -- & 49.2 & 57.8 & -- & -- \\
$\ge$T3, \% & 5.2 & 16.3 & 12.4 & -- & -- & -- & 8.2 & 16.4 & -- & -- \\
Node positive, \% & 47.6 & 50.3 & 40.3 & -- & -- & -- & 49.1 & 51.7 & -- & -- \\
Post-menopausal, \% & 78.0 & 48.0 & 53.8 & 28.1 & -- & 30.8 & -- & -- & -- & -- \\
\bottomrule
\end{tabular}}
\end{table}

\begin{figure}[htbp]\centering
\adjustbox{max width=\linewidth}{\pgfplotstableread{\gdata{grade_distribution}}\tblgrade
\begin{tikzpicture}
\begin{axis}[bcbase, xbar stacked, width=0.9\linewidth, height=6.2cm, bar width=8pt,
  xmin=0, xmax=100, xlabel={Patients (\%)}, xmajorgrids,
  ytick=data, yticklabels from table={\tblgrade}{label}, y dir=normal,
  enlarge y limits=0.08, legend columns=3, legend style={at={(0.5,1.02)}, anchor=south},
  nodes near coords={\pgfmathprintnumber[fixed,precision=0]{\pgfplotspointmeta}},
  nodes near coords style={font=\sffamily\tiny, text=white}, point meta=explicit,
  every node near coord/.append style={anchor=center}]
\addplot[fill=gI, draw=white] table[x=g1, y=y, meta=g1] {\tblgrade};
\addplot[fill=gII, draw=white] table[x=g2, y=y, meta=g2] {\tblgrade};
\addplot[fill=gIII, draw=white] table[x=g3, y=y, meta=g3] {\tblgrade};
\legend{Grade 1, Grade 2, Grade 3}
\draw[inkb, dashed] (axis cs:0,1.5) -- (axis cs:100,1.5);
\end{axis}
\end{tikzpicture}
}
\caption{Distribution of Nottingham grade by cohort (\%), with the published Saudi series below the
dashed line ($^\dagger$\citealp{kussaibi2021saudi}; $^\ddagger$\citealp{durrani2025healthcare}).}\label{fig:grade}
\end{figure}

\subsection{Associations between grade and routine variables}
Of the \NTests{} pre-specified associations, \NSigTests{} remained significant after adjustment for
multiple testing (Table~\ref{tab:assoc}, Figure~\ref{fig:assoc}). In METABRIC, high grade was most
strongly linked to ER-negative ($\tau_b$ \TauER) and PR-negative ($\tau_b$ \TauPR) disease, followed
by HER2 positivity (\TauHER), larger tumours (\TauT), nodal involvement (\TauN) and younger age
($\rho$ \RhoAge). Histological type showed a moderate association ($V$ = \VHist). All directions were
the same in MSK-2018. In the multivariable proportional-odds model (\NOrdinal{} complete cases;
Table~\ref{tab:shap}), ER positivity (odds ratio for higher grade \OrER, 95\% CI \OrERlo--\OrERhi) and
lobular histology (\OrILC, \OrILClo--\OrILChi) remained the strongest independent factors.

\begin{table}[htbp]\centering\small
\caption{Scale-appropriate associations with bootstrap 95\% CIs and Benjamini--Hochberg-adjusted
$p$ values. Binary variables are coded 1 = present/positive, so a negative $\tau_b$ for ER means
ER-positive tumours tend to be lower grade.}\label{tab:assoc}
\tabfit{\begin{tabular}{llrrrr}
\toprule
Pair & Measure & \multicolumn{2}{c}{METABRIC} & \multicolumn{2}{c}{MSK-2018} \\
\cmidrule(lr){3-4}\cmidrule(lr){5-6}
 & & Effect (95\% CI) & $p_\mathrm{BH}$ & Effect (95\% CI) & $p_\mathrm{BH}$ \\
\midrule
Grade $\times$ ER+ & $\tau_b$ & -0.373 (-0.404, -0.342) & $<$0.001 & -0.211 (-0.241, -0.178) & $<$0.001 \\
Grade $\times$ PR+ & $\tau_b$ & -0.310 (-0.350, -0.271) & $<$0.001 & -0.203 (-0.245, -0.165) & $<$0.001 \\
Grade $\times$ HER2+ & $\tau_b$ & +0.194 (+0.155, +0.229) & $<$0.001 & +0.143 (+0.106, +0.181) & $<$0.001 \\
Grade $\times$ T category & $\tau_b$ & +0.153 (+0.115, +0.191) & $<$0.001 & +0.210 (+0.166, +0.257) & $<$0.001 \\
Grade $\times$ Node+ & $\tau_b$ & +0.137 (+0.097, +0.178) & $<$0.001 & +0.172 (+0.117, +0.221) & $<$0.001 \\
Grade $\times$ Post-menopausal & $\tau_b$ & -0.104 (-0.148, -0.061) & $<$0.001 & -0.043 (-0.091, +0.004) & 0.083 \\
Age $\times$ Grade & $\rho$ & -0.118 (-0.158, -0.078); $\varepsilon^2$=0.016 & $<$0.001 & -0.123 (-0.171, -0.074); $\varepsilon^2$=0.016 & $<$0.001 \\
Grade $\times$ Histology & $V$ & +0.177 (+0.145, +0.217) & $<$0.001 & +0.108 (+0.076, +0.157) & $<$0.001 \\
Histology $\times$ ER+ & $V$ & +0.139 (+0.109, +0.174) & $<$0.001 & +0.205 (+0.158, +0.263) & $<$0.001 \\
T category $\times$ Node+ & $\tau_b$ & +0.218 (+0.179, +0.256) & $<$0.001 & +0.374 (+0.330, +0.417) & $<$0.001 \\
Age $\times$ ER+ & $\rho$ & +0.267 (+0.227, +0.304); $\varepsilon^2$=0.071 & $<$0.001 & +0.017 (-0.033, +0.065); $\varepsilon^2$=0.000 & 0.488 \\
\bottomrule
\end{tabular}}
\end{table}

\begin{figure}[htbp]\centering
\adjustbox{max width=\linewidth}{\pgfplotstableread{\gdata{associations}}\tblassoc
\begin{tikzpicture}
\begin{axis}[bcbase, width=0.62\linewidth, height=6cm, xlabel={Effect size (bootstrap 95\% CI)},
  ytick=data, yticklabels from table={\tblassoc}{label}, xmajorgrids, enlarge y limits=0.08,
  legend columns=2, legend style={at={(0.5,1.02)}, anchor=south}]
\addplot[dotci=cA] table[x=met, y expr=\thisrow{y}+0.15, x error minus=met_em, x error plus=met_ep] {\tblassoc};
\addplot[dotci=cB] table[x=msk, y expr=\thisrow{y}-0.15, x error minus=msk_em, x error plus=msk_ep] {\tblassoc};
\legend{METABRIC, MSK-2018}
\draw[inkb] (axis cs:0,\pgfkeysvalueof{/pgfplots/ymin}) -- (axis cs:0,\pgfkeysvalueof{/pgfplots/ymax});
\end{axis}
\end{tikzpicture}
}
\caption{Associations involving grade in the two large cohorts (effect and bootstrap 95\% CI).}\label{fig:assoc}
\end{figure}

\subsection{Internal validation: flexible models did not beat regression}
All learned models performed well above the baseline, and they performed almost identically
(Table~\ref{tab:internal}, Figure~\ref{fig:ladder}). Multinomial LR (macro-AUC \MlrIntAuc) and the
proportional-odds model (\OrdIntAuc) matched XGBoost (\XgbIntAuc) and the other ensemble methods.
The paired difference between XGBoost and LR was \XgbDeltaInt{} (95\% CI \XgbDeltaIntLo{} to
\XgbDeltaIntHi), and no model exceeded LR by more than \MaxDeltaIntHi{} even at the upper confidence
limit. $k$-nearest neighbours performed worst (\KnnIntAuc). Results were stable across repeats
(standard deviation at most \MaxRepSd). Grade~3 was the easiest class to identify (AUC
\BestIntGThreeAuc{} for \BestIntModel) and grade~2 the hardest. Macro-F1 was \MlrIntFOne{} for LR
compared with \MajIntFOne{} for the baseline. Multiple imputation changed macro-AUC by at most
\MIDiffMax{} (Supplementary Table~\ref{tab:mi}), and a complete-case analysis (\NComplete{} patients)
gave \CCMlrAuc{} for LR. The sample-size criteria were \SSMet: \SSRequired{} patients would be needed,
against \SSAvailable{} available. The limiting factor was the grade~1 versus grade~2 comparison,
which carries little predictive signal (Supplementary Table~\ref{tab:ss}).

\begin{table}[htbp]\centering\small
\caption{Internal validation in METABRIC (nested, \Repeats$\times$ repeated stratified \Folds-fold
cross-validation; full feature set). Point estimates are means over repeats; 95\% CIs are
patient-level bootstraps (\Boot{} resamples) of repeat-averaged out-of-fold predictions. $\Delta$ is
the paired-bootstrap difference in macro-AUC versus multinomial LR. QWK = quadratic-weighted
$\kappa$.}\label{tab:internal}
\tabfit{\begin{tabular}{lrrrrrrrr}
\toprule
Model & Macro-AUC (95\% CI) & $\Delta$ vs LR (95\% CI) & AUC G1 & AUC G2 & AUC G3 & Macro-F1 (95\% CI) & QWK & Brier \\
\midrule
Majority class & 0.499 (0.478--0.514) & -0.228 (-0.255, -0.203) & 0.497 & 0.500 & 0.499 & 0.223 (0.217--0.229) & 0.000 & 0.573 \\
Multinomial LR & 0.725 (0.707--0.742) & ref. & 0.741 & 0.681 & 0.752 & 0.465 (0.445--0.486) & 0.356 & 0.490 \\
Ordinal LR (prop. odds) & 0.725 (0.709--0.743) & +0.001 (-0.004, +0.006) & 0.744 & 0.679 & 0.752 & 0.458 (0.435--0.476) & 0.366 & 0.491 \\
k-NN & 0.715 (0.699--0.735) & -0.008 (-0.016, +0.001) & 0.729 & 0.676 & 0.741 & 0.452 (0.427--0.468) & 0.315 & 0.496 \\
Random forest & 0.724 (0.708--0.743) & +0.001 (-0.007, +0.009) & 0.739 & 0.682 & 0.751 & 0.449 (0.424--0.463) & 0.344 & 0.492 \\
XGBoost & 0.723 (0.708--0.743) & +0.001 (-0.006, +0.008) & 0.737 & 0.678 & 0.753 & 0.458 (0.440--0.480) & 0.344 & 0.492 \\
LightGBM & 0.719 (0.705--0.741) & -0.001 (-0.010, +0.007) & 0.728 & 0.678 & 0.752 & 0.457 (0.443--0.484) & 0.338 & 0.496 \\
\bottomrule
\end{tabular}}
\end{table}

\begin{figure*}[htbp]\centering
\resizebox{\linewidth}{!}{\pgfplotstableread{\gdata{model_ladder}}\tblladder
\begin{tikzpicture}
\begin{groupplot}[group style={group size=4 by 1, horizontal sep=4mm, yticklabels at=edge left},
  bcbase, width=0.31\linewidth, height=5.2cm, xmin=0.45, xmax=0.88, xmajorgrids,
  ytick=data, yticklabels from table={\tblladder}{label}, enlarge y limits=0.1,
  xtick={0.5,0.6,0.7,0.8}, title style={font=\sffamily\scriptsize\bfseries, yshift=-1mm}]
\nextgroupplot[title={Internal (METABRIC CV)}]
\addplot[dotci=cA] table[x=int, y=y, x error minus=int_em, x error plus=int_ep] {\tblladder};
\draw[inkb, dashed] (axis cs:0.5,-1) -- (axis cs:0.5,10);
\nextgroupplot[title={MSK-2018 (full)}]
\addplot[dotci=cB] table[x=msk, y=y, x error minus=msk_em, x error plus=msk_ep] {\tblladder};
\draw[inkb, dashed] (axis cs:0.5,-1) -- (axis cs:0.5,10);
\nextgroupplot[title={Duke MRI (full)}]
\addplot[dotci=cD] table[x=duke, y=y, x error minus=duke_em, x error plus=duke_ep] {\tblladder};
\draw[inkb, dashed] (axis cs:0.5,-1) -- (axis cs:0.5,10);
\nextgroupplot[title={MBCproject (core)}]
\addplot[dotci=cC] table[x=mbcp, y=y, x error minus=mbcp_em, x error plus=mbcp_ep] {\tblladder};
\draw[inkb, dashed] (axis cs:0.5,-1) -- (axis cs:0.5,10);
\end{groupplot}
\node[font=\sffamily\scriptsize, anchor=north] at ($(group c2r1.south east)!0.5!(group c3r1.south west)+(0,-6mm)$) {Macro-averaged one-vs-rest AUC (95\% CI); dashed = chance};
\end{tikzpicture}
}
\caption{Macro-averaged AUC (95\% CI) of the model ladder internally and in three external cohorts;
dashed line = chance.}\label{fig:ladder}
\end{figure*}

\subsection{External and leave-one-dataset-out validation}
When applied unchanged to new cohorts with the full feature set, the models reached macro-AUCs of
\FullMskMin--\FullMskMax{} in MSK-2018 (LR \MlrExtAuc, 95\% CI \MlrMskAucLo--\MlrMskAucHi) and
\FullDukeMin--\FullDukeMax{} in Duke (LR \MlrDukeAuc, \MlrDukeAucLo--\MlrDukeAucHi; \NDukeExt{}
patients) (Table~\ref{tab:external}). This is close to the internal estimate, so discrimination
transported well. With the core feature set, which lacks tumour size and nodal status, MSK-2018
discrimination fell to \CoreMskMin--\CoreMskMax. The smaller cohorts reached up to \CoreExtAucMax,
with wide confidence intervals. Leave-one-dataset-out estimates ranged from \LodoAucMin{} to
\LodoAucMax{} (Supplementary Table~\ref{tab:lodo}). Of \NExtComparisons{} paired external
comparisons between a flexible model and LR, \NExtSigBetter{} favoured the flexible model and
\NExtSigWorse{} favoured LR.

\begin{table}[htbp]\centering\small
\caption{External validation of models developed in METABRIC and applied without refitting.
Calibration slope is for P(grade 3).}\label{tab:external}
\tabfit{\begin{tabular}{llrrrrr}
\toprule
Target (features) & Model & $n$ & Macro-AUC (95\% CI) & G3-AUC (95\% CI) & Macro-F1 & Cal. slope (G3) \\
\midrule
MSK-2018 (full) & Majority class & 1573 & 0.500 (0.500--0.500) & 0.500 (0.500--0.500) & 0.269 & -- \\
MSK-2018 (full) & Multinomial LR & 1573 & 0.710 (0.685--0.734) & 0.727 (0.701--0.753) & 0.372 & 1.03 \\
MSK-2018 (full) & Ordinal LR (prop. odds) & 1573 & 0.713 (0.688--0.737) & 0.725 (0.699--0.750) & 0.370 & 1.02 \\
MSK-2018 (full) & Random forest & 1573 & 0.704 (0.678--0.728) & 0.725 (0.699--0.751) & 0.350 & 1.11 \\
MSK-2018 (full) & XGBoost & 1573 & 0.711 (0.686--0.735) & 0.726 (0.701--0.753) & 0.373 & 1.03 \\
MSK-2018 (full) & LightGBM & 1573 & 0.705 (0.681--0.730) & 0.720 (0.693--0.746) & 0.388 & 0.82 \\
\addlinespace
Duke MRI (full) & Majority class & 637 & 0.500 (0.500--0.500) & 0.500 (0.500--0.500) & 0.164 & -- \\
Duke MRI (full) & Multinomial LR & 637 & 0.709 (0.681--0.737) & 0.779 (0.739--0.816) & 0.494 & 0.96 \\
Duke MRI (full) & Ordinal LR (prop. odds) & 637 & 0.705 (0.678--0.731) & 0.778 (0.738--0.814) & 0.487 & 0.92 \\
Duke MRI (full) & Random forest & 637 & 0.717 (0.689--0.744) & 0.789 (0.752--0.826) & 0.467 & 1.01 \\
Duke MRI (full) & XGBoost & 637 & 0.720 (0.691--0.748) & 0.789 (0.750--0.826) & 0.486 & 0.97 \\
Duke MRI (full) & LightGBM & 637 & 0.720 (0.691--0.748) & 0.786 (0.749--0.824) & 0.484 & 0.88 \\
\addlinespace
MSK-2018 (core) & Majority class & 1573 & 0.500 (0.500--0.500) & 0.500 (0.500--0.500) & 0.269 & -- \\
MSK-2018 (core) & Multinomial LR & 1573 & 0.667 (0.642--0.692) & 0.685 (0.658--0.711) & 0.349 & 0.90 \\
MSK-2018 (core) & Ordinal LR (prop. odds) & 1573 & 0.670 (0.645--0.695) & 0.682 (0.655--0.709) & 0.357 & 0.89 \\
MSK-2018 (core) & Random forest & 1573 & 0.665 (0.639--0.690) & 0.690 (0.663--0.717) & 0.357 & 0.98 \\
MSK-2018 (core) & XGBoost & 1573 & 0.658 (0.632--0.684) & 0.680 (0.652--0.707) & 0.357 & 0.88 \\
MSK-2018 (core) & LightGBM & 1573 & 0.660 (0.633--0.686) & 0.681 (0.654--0.708) & 0.357 & 0.72 \\
\addlinespace
MBCproject (core) & Majority class & 207 & 0.500 (0.500--0.500) & 0.500 (0.500--0.500) & 0.182 & -- \\
MBCproject (core) & Multinomial LR & 207 & 0.701 (0.647--0.754) & 0.771 (0.703--0.832) & 0.486 & 1.16 \\
MBCproject (core) & Ordinal LR (prop. odds) & 207 & 0.715 (0.665--0.761) & 0.762 (0.692--0.829) & 0.513 & 1.08 \\
MBCproject (core) & Random forest & 207 & 0.693 (0.640--0.744) & 0.753 (0.686--0.819) & 0.440 & 1.12 \\
MBCproject (core) & XGBoost & 207 & 0.697 (0.642--0.750) & 0.758 (0.689--0.822) & 0.455 & 1.01 \\
MBCproject (core) & LightGBM & 207 & 0.703 (0.649--0.757) & 0.760 (0.693--0.826) & 0.485 & 0.93 \\
\addlinespace
Sanger (core) & Majority class & 100 & 0.500 (0.500--0.500) & 0.500 (0.500--0.500) & 0.225 & -- \\
Sanger (core) & Multinomial LR & 100 & 0.797 (0.712--0.872) & 0.806 (0.714--0.882) & 0.570 & 1.33 \\
Sanger (core) & Ordinal LR (prop. odds) & 100 & 0.763 (0.679--0.842) & 0.801 (0.704--0.882) & 0.495 & 1.25 \\
Sanger (core) & Random forest & 100 & 0.764 (0.680--0.840) & 0.772 (0.676--0.856) & 0.441 & 1.27 \\
Sanger (core) & XGBoost & 100 & 0.766 (0.678--0.847) & 0.774 (0.677--0.859) & 0.479 & 1.20 \\
Sanger (core) & LightGBM & 100 & 0.767 (0.684--0.845) & 0.774 (0.675--0.862) & 0.486 & 1.03 \\
\addlinespace
Broad (MX/VN) (core) & Majority class & 59 & 0.500 (0.500--0.500) & 0.500 (0.500--0.500) & 0.204 & -- \\
Broad (MX/VN) (core) & Multinomial LR & 59 & 0.784 (0.691--0.869) & 0.737 (0.599--0.863) & 0.663 & 0.90 \\
Broad (MX/VN) (core) & Ordinal LR (prop. odds) & 59 & 0.777 (0.687--0.863) & 0.738 (0.604--0.866) & 0.675 & 0.86 \\
Broad (MX/VN) (core) & Random forest & 59 & 0.795 (0.678--0.891) & 0.793 (0.653--0.910) & 0.403 & 1.30 \\
Broad (MX/VN) (core) & XGBoost & 59 & 0.787 (0.668--0.882) & 0.773 (0.634--0.888) & 0.434 & 1.04 \\
Broad (MX/VN) (core) & LightGBM & 59 & 0.785 (0.665--0.887) & 0.760 (0.616--0.884) & 0.562 & 0.96 \\
\addlinespace
Duke MRI (core) & Majority class & 637 & 0.500 (0.500--0.500) & 0.500 (0.500--0.500) & 0.164 & -- \\
Duke MRI (core) & Multinomial LR & 637 & 0.702 (0.675--0.729) & 0.783 (0.743--0.820) & 0.434 & 1.01 \\
Duke MRI (core) & Ordinal LR (prop. odds) & 637 & 0.705 (0.678--0.732) & 0.782 (0.743--0.819) & 0.476 & 0.98 \\
Duke MRI (core) & Random forest & 637 & 0.715 (0.686--0.743) & 0.782 (0.744--0.821) & 0.485 & 1.02 \\
Duke MRI (core) & XGBoost & 637 & 0.714 (0.684--0.743) & 0.784 (0.746--0.822) & 0.476 & 0.97 \\
Duke MRI (core) & LightGBM & 637 & 0.712 (0.683--0.739) & 0.779 (0.740--0.817) & 0.487 & 0.88 \\
\bottomrule
\end{tabular}}
\end{table}

\subsection{Calibration, updating and clinical utility}
Predictions were well calibrated internally (calibration slope \CalSlopeInt{} for LR and
\CalSlopeIntXgb{} for XGBoost; Supplementary Table~\ref{tab:cal}). In the external cohorts the slope
stayed close to one (MSK-2018 \CalSlopeMskRaw; Duke \CalSlopeDukeRaw), which means the predictor
effects transported. The average level of the predictions did not. In MSK-2018, \PrevMsk{} of
patients had grade~3 but the mean predicted probability was only \PredMsk{} (calibration intercept
\CalIntMskRaw), so the model under-predicted grade~3. In Duke the opposite occurred: \PrevDuke{}
observed against \PredDuke{} predicted (intercept \CalIntDukeRaw). Re-estimating only the intercept
in each cohort removed this miscalibration without changing discrimination (Figure~\ref{fig:cal};
per-grade results in Supplementary Table~\ref{tab:oe}). Decision curves showed a net benefit over
treating all or no patients as grade~3 for threshold probabilities above about \DcaLoInt{}
internally and \DcaLoExt{} in MSK-2018, where the high prevalence of grade~3 makes ``treat all'' a
strong alternative, and between \DcaLoDuke{} and \DcaHiDuke{} in Duke.

\begin{figure*}[htbp]\centering
\adjustbox{max width=\linewidth}{\begin{tikzpicture}
\begin{groupplot}[group style={group size=2 by 1, horizontal sep=16mm}, bcbase, width=0.47\linewidth, height=6cm,
  grid=major, title style={font=\sffamily\scriptsize\bfseries}]
\nextgroupplot[title={Calibration of P(grade 3), multinomial LR}, xmin=0, xmax=1, ymin=0, ymax=1,
  xlabel={Predicted probability}, ylabel={Observed proportion (LOWESS)},
  legend style={at={(0.02,0.98)}, anchor=north west}]
\addplot[cA, very thick] table[x=p, y=obs] {\gdata{cal_int}};
\addplot[cB, very thick] table[x=p, y=obs] {\gdata{cal_msk}};
\addplot[cB, very thick, dashed] table[x=p, y=obs] {\gdata{cal_mskrec}};
\addplot[cD, very thick] table[x=p, y=obs] {\gdata{cal_duke}};
\addplot[cD, very thick, dashed] table[x=p, y=obs] {\gdata{cal_dukerec}};
\addplot[inkb, densely dotted, domain=0:1] {x};
\legend{METABRIC (internal), MSK-2018, MSK-2018 updated, Duke MRI, Duke updated}
\nextgroupplot[title={Decision curves for grade 3}, xmin=0.05, xmax=0.85, ymin=-0.05, ymax=0.7,
  xlabel={Threshold probability}, ylabel={Net benefit}, legend style={at={(0.98,0.98)}, anchor=north east}]
\addplot[cA, very thick] table[x=threshold, y=model_int] {\gdata{dca}};
\addplot[cB, very thick] table[x=threshold, y=model_msk] {\gdata{dca}};
\addplot[cD, very thick] table[x=threshold, y=model_duke] {\gdata{dca}};
\addplot[cA, thin, dashed, forget plot] table[x=threshold, y=all_int] {\gdata{dca}};
\addplot[cB, thin, dashed, forget plot] table[x=threshold, y=all_msk] {\gdata{dca}};
\addplot[cD, thin, dashed, forget plot] table[x=threshold, y=all_duke] {\gdata{dca}};
\addplot[inkb, thin, domain=0.05:0.85] {0};
\legend{METABRIC model, MSK-2018 model, Duke model, Treat none}
\end{groupplot}
\end{tikzpicture}
}
\caption{Left: LOWESS calibration curves for P(grade 3) from multinomial LR, internally, externally
(solid) and after a cross-fitted intercept update (dashed). Right: decision curves (solid = model;
thin dashed = treat all, same colour; treat none = 0).}\label{fig:cal}
\end{figure*}

\subsection{Interpretability and clinical plausibility}
The SHAP analysis ranked \ShapTop{} as the most important feature and \ShapSecond{} as the second
(Figure~\ref{fig:shap}). \ShapAgree{} pre-specified directions were confirmed, with bootstrap
intervals excluding zero (Table~\ref{tab:shap}), and all agreed in sign with the proportional-odds
odds ratios. The model therefore uses the variables in the way a pathologist would expect.

\begin{figure}[htbp]\centering
\adjustbox{max width=\linewidth}{\pgfplotstableread{\gdata{shap}}\tblshap
\begin{tikzpicture}
\begin{groupplot}[group style={group size=2 by 1, horizontal sep=4mm, yticklabels at=edge left}, bcbase,
  width=0.42\linewidth, height=6cm, ytick=data, yticklabels from table={\tblshap}{label}, enlarge y limits=0.07,
  xmajorgrids, title style={font=\sffamily\scriptsize\bfseries}]
\nextgroupplot[title={Importance: mean $|$SHAP$|$ (grade 3)}, xbar, bar width=6pt, xmin=0, xlabel={log-odds}]
\addplot[fill=cA, draw=none] table[x=meanabs, y=y] {\tblshap};
\nextgroupplot[title={Direction: Spearman $\rho$ (value, SHAP)}, xmin=-1, xmax=1, xlabel={$\rho$ (bootstrap 95\% CI)}]
\addplot[dotci=cB] table[x=rho, y=y, x error minus=rho_em, x error plus=rho_ep] {\tblshap};
\draw[inkb] (axis cs:0,-1) -- (axis cs:0,20);
\end{groupplot}
\end{tikzpicture}
}
\caption{Global importance (mean $|$SHAP$|$) and direction (Spearman $\rho$ between feature value and
SHAP value, bootstrap 95\% CI) for the grade-3 output of XGBoost in METABRIC. Beeswarm in
Supplementary Figure~\ref{fig:beeswarm}.}\label{fig:shap}
\end{figure}

\begin{table}[htbp]\centering\small
\caption{Clinical-plausibility check. ``Expected'' was pre-specified. Ordinal odds ratios are for
higher grade from a complete-case proportional-odds model (bootstrap 95\% CI).}\label{tab:shap}
\tabfit{\begin{tabular}{lrrlrl}
\toprule
Feature & mean$|$SHAP$|$ & Direction $\rho$ (95\% CI) & Expected & Agrees & Ordinal OR (95\% CI) \\
\midrule
ER positive & 0.330 & -0.76 (-0.78, -0.74) & $\downarrow$ & yes & 0.17 (0.12--0.24) \\
T category & 0.158 & +0.82 (+0.81, +0.83) & $\uparrow$ & yes & 1.61 (1.31--2.00) \\
PR positive & 0.135 & -0.87 (-0.87, -0.86) & $\downarrow$ & yes & 0.57 (0.46--0.72) \\
Node positive & 0.086 & +0.87 (+0.86, +0.87) & $\uparrow$ & yes & 1.29 (1.08--1.62) \\
HER2 positive & 0.065 & +0.57 (+0.54, +0.60) & $\uparrow$ & yes & 1.83 (1.28--2.69) \\
Mixed (vs IDC) & 0.064 & -0.55 (-0.57, -0.52) & -- & -- & 0.53 (0.38--0.74) \\
Age (per 10 y) & 0.061 & -0.25 (-0.30, -0.21) & $\downarrow$ & yes & 0.93 (0.83--1.03) \\
ILC (vs IDC) & 0.053 & -0.47 (-0.49, -0.43) & $\downarrow$ & yes & 0.41 (0.28--0.57) \\
Other type (vs IDC) & 0.034 & -0.35 (-0.38, -0.31) & -- & -- & 0.26 (0.16--0.43) \\
Post-menopausal & 0.008 & -0.65 (-0.68, -0.61) & -- & -- & 0.96 (0.68--1.37) \\
\bottomrule
\end{tabular}}
\end{table}

\subsection{Leakage audit}
Every shortcut inflated apparent performance (Table~\ref{tab:leak}, Figure~\ref{fig:leak}). Adding the
Nottingham Prognostic Index raised the macro-AUC of XGBoost from \NpiHonest{} to \NpiLeaky. Adding the
Nottingham component scores raised LR from \CompHonest{} to \CompLeaky{} in Sanger and from
\CompHonestDuke{} to \CompLeakyDuke{} in Duke. Imputing and oversampling before splitting raised
$k$-nearest neighbours from \SmoteHonestKnn{} to \SmoteLeakyKnn. Splitting by record rather than by
patient had a smaller effect (\RecHonestKnn{} to \RecLeakyKnn{} for 1-nearest neighbour) because few
MSK-2018 patients had more than one record. The largest inflation, \LeakMaxInfl{} in macro-AUC, is
larger than the entire gap between the baseline and the best correctly validated model.

\begin{table}[htbp]\centering\small
\caption{Leakage audit: identical models evaluated with the correct and a leaky protocol (5-fold
cross-validation).}\label{tab:leak}
\tabfit{\begin{tabular}{p{6.2cm}lrrr}
\toprule
Scenario & Model & Honest AUC & Leaky AUC & Inflation \\
\midrule
Target leakage: NPI as predictor (METABRIC) & Multinomial LR & 0.726 & 0.949 & +0.223 \\
Target leakage: NPI as predictor (METABRIC) & XGBoost & 0.721 & 0.972 & +0.251 \\
SMOTE + imputation before split (METABRIC) & k-NN & 0.652 & 0.806 & +0.154 \\
SMOTE + imputation before split (METABRIC) & XGBoost & 0.669 & 0.850 & +0.181 \\
Grade components as predictors (Sanger) & Multinomial LR & 0.696 & 0.962 & +0.266 \\
Grade components as predictors (Duke MRI) & Multinomial LR & 0.698 & 0.995 & +0.297 \\
Record-level split, 1718 records / 1573 patients (MSK) & k-NN & 0.544 & 0.579 & +0.035 \\
Record-level split, 1718 records / 1573 patients (MSK) & Random forest & 0.621 & 0.653 & +0.033 \\
\bottomrule
\end{tabular}}
\end{table}

\begin{figure}[htbp]\centering
\adjustbox{max width=\linewidth}{\pgfplotstableread{\gdata{leakage}}\tblleak
\begin{tikzpicture}
\begin{axis}[bcbase, width=0.55\linewidth, height=6cm, xlabel={Cross-validated macro-AUC},
  ytick=data, yticklabels from table={\tblleak}{label}, xmajorgrids, enlarge y limits=0.08,
  legend columns=2, legend style={at={(0.5,1.02)}, anchor=south}]
\addplot[only marks, mark=*, mark size=1.8pt, cA, error bars/.cd, x dir=plus, x explicit, error mark=none,
  error bar style={gridc, line width=2pt}]
  table[x=honest, y=y, x error expr=\thisrow{leaky}-\thisrow{honest}] {\tblleak};
\addplot[only marks, mark=*, mark size=1.8pt, cB] table[x=leaky, y=y] {\tblleak};
\legend{Correct protocol, Leaky protocol}
\end{axis}
\end{tikzpicture}
}
\caption{Apparent macro-AUC under the correct (blue) and leaky (orange) protocols.}\label{fig:leak}
\end{figure}

\subsection{Imaging extension: MRI radiomics in Duke}
In the \NDukeRad{} graded Duke patients, the clinical model reached a cross-validated macro-AUC of
\DukeClinAuc{} (95\% CI \DukeClinAucLo--\DukeClinAucHi). Radiomics alone reached \DukeRadAuc{}
(\DukeRadAucLo--\DukeRadAucHi), and the combined model \DukeFusAuc{} (\DukeFusAucLo--\DukeFusAucHi)
(Table~\ref{tab:duke}, Figure~\ref{fig:duke}). Adding radiomics to the clinical variables therefore
lowered performance, by \DukeFusDelta{} (\DukeFusDeltaLo{} to \DukeFusDeltaHi).

\begin{table}[htbp]\centering\small
\caption{Duke-Breast-Cancer-MRI: clinical, radiomics and fusion models (nested
\RadRepeats$\times$\Folds-fold cross-validation; $L_2$-penalised multinomial LR; paired bootstrap
$\Delta$ versus the clinical model).}\label{tab:duke}
\tabfit{\begin{tabular}{lrrrrr}
\toprule
Model (L2 logistic) & Features & Macro-AUC (95\% CI) & G3-AUC (95\% CI) & $\Delta$ macro-AUC vs clinical (95\% CI) & Brier \\
\midrule
Clinical (full set) & 10 & 0.694 (0.663--0.727) & 0.769 (0.731--0.811) & ref. & 0.551 \\
DCE-MRI radiomics (529 features) & 529 & 0.608 (0.575--0.645) & 0.667 (0.626--0.718) & -0.085 (-0.126, -0.045) & 0.600 \\
Clinical + radiomics & 539 & 0.648 (0.618--0.685) & 0.731 (0.693--0.777) & -0.045 (-0.077, -0.011) & 0.582 \\
\bottomrule
\end{tabular}}
\end{table}

\begin{figure}[htbp]\centering
\adjustbox{max width=\linewidth}{\pgfplotstableread{\gdata{duke_radiomics}}\tbldk
\begin{tikzpicture}
\begin{axis}[bcbase, width=0.55\linewidth, height=4.2cm, xmin=0.45, xmax=0.9, xmajorgrids,
  ytick=data, yticklabels from table={\tbldk}{label}, enlarge y limits=0.25, xlabel={AUC (95\% CI), nested CV in Duke},
  legend columns=2, legend style={at={(0.5,1.02)}, anchor=south}]
\addplot[dotci=cA] table[x=macro, y expr=\thisrow{y}+0.12, x error minus=macro_em, x error plus=macro_ep] {\tbldk};
\addplot[dotci=cD] table[x=g3, y expr=\thisrow{y}-0.12, x error minus=g3_em, x error plus=g3_ep] {\tbldk};
\legend{Macro-AUC (3 grades), Grade 3 vs 1--2}
\draw[inkb, dashed] (axis cs:0.5,-1) -- (axis cs:0.5,5);
\end{axis}
\end{tikzpicture}
}
\caption{Duke: discrimination of clinical, radiomics and fusion models (95\% CI); dashed = chance.}\label{fig:duke}
\end{figure}

\subsection{Transport to Saudi case-mix}
When METABRIC was re-weighted to the Dammam receptor, histology and age profile (effective sample
size \EssMet{} of \NRakedMet), its grade distribution moved to \RwGOneMet/\RwGTwoMet/\RwGThreeMet\%
(grade~3 95\% CI \RwGThreeMetLo--\RwGThreeMetHi\%), compared with \ObsGThreeMet\% grade~3 observed
in METABRIC and 36.7\% published for Dammam (Table~\ref{tab:transport}, Figure~\ref{fig:transport}).
In other words, the Saudi case-mix predicts \emph{more} grade~3, not less. The same pattern appeared
for MSK-2018 (\RwGThreeMsk\%) and Duke (\RwGThreeDuke\%), for every assumed share of patients aged
$\le$50 years from \AgeSensLo\% to \AgeSensHi\% (grade~3 \AgeSensMin--\AgeSensMax\%; Supplementary
Table~\ref{tab:agesens}), and for the Riyadh series (re-weighted METABRIC grade~3 \RwGThreeMetB\%,
\RwGThreeMetBLo--\RwGThreeMetBHi\%, against \SaudiGThreeB\% published). Across \PoolK{} Saudi series
(\PoolN{} patients), the pooled published grade-3 share was \PoolGThree\% (\PoolGThreeLo--\PoolGThreeHi\%);
including the Najran series~\citep{alamri2024najran} gave \PoolGThreeAll\%
(\PoolGThreeAllLo--\PoolGThreeAllHi\%). Under the Dammam case-mix, LR discrimination rose from
\AucUnwMet{} to \AucWMet{}, but the mean predicted probability of grade~3 was \PredGThreeMet\%,
about \OverPredPP{} percentage points above the published Dammam figure.

\begin{table}[htbp]\centering\small
\caption{Re-weighting to published Saudi case-mix. ESS = effective sample size of the raking weights
(first row of each block: patients with complete raking variables). Grade-3 CI: bootstrap with
re-raking. Arrows: unweighted$\to$Saudi-weighted LR discrimination. $^\dagger$Dammam
series~\citep{kussaibi2021saudi}; $^\ddagger$Riyadh series~\citep{durrani2025healthcare}
(tertiary referral centre). Pooled: DerSimonian--Laird random effects across \PoolK{} Saudi
series reporting grade~\citep{kussaibi2021saudi,durrani2025healthcare,albasri2014madinah}; the
Madinah series reports grade for invasive ductal carcinoma only and no receptor status, so it
contributes to the pooled estimate but not to raking.}\label{tab:transport}
\tabfit{\begin{tabular}{lllrrrrrr}
\toprule
Anchor & Cohort & Quantity & G1 & G2 & G3 (95\% CI) & ESS & Macro-AUC & G3-AUC \\
\midrule
Dammam (Kussaibi 2021)$^\dagger$ & METABRIC & observed & 8.9 & 40.5 & 50.6 & 1861 & & \\
 & & re-weighted & 6.2 & 32.9 & 60.9 (58.5--63.0) & 1088 & & \\
 & & LR predicted (weighted) & 6.7 & 34.2 & 59.1 & & 0.727$\to$0.790 & 0.759$\to$0.811 \\
 & & published & 8.7 & 54.5 & 36.7 & & & \\
\addlinespace
Dammam (Kussaibi 2021)$^\dagger$ & MSK-2018 & observed & 5.7 & 26.7 & 67.6 & 1496 & & \\
 & & re-weighted & 3.4 & 19.2 & 77.3 (75.1--79.9) & 846 & & \\
 & & LR predicted (weighted) & 5.6 & 33.1 & 61.3 & & 0.715$\to$0.766 & 0.734$\to$0.765 \\
 & & published & 8.7 & 54.5 & 36.7 & & & \\
\addlinespace
Dammam (Kussaibi 2021)$^\dagger$ & Duke MRI & observed & 17.7 & 49.8 & 32.5 & 637 & & \\
 & & re-weighted & 13.2 & 44.7 & 42.1 (37.7--45.9) & 510 & & \\
 & & LR predicted (weighted) & 5.8 & 32.8 & 61.4 & & 0.709$\to$0.717 & 0.779$\to$0.755 \\
 & & published & 8.7 & 54.5 & 36.7 & & & \\
\addlinespace
Riyadh (Durrani 2025)$^\ddagger$ & METABRIC & observed & 8.9 & 40.5 & 50.6 & 1861 & & \\
 & & re-weighted & 8.2 & 34.5 & 57.3 (54.6--60.0) & 852 & & \\
 & & LR predicted (weighted) & 8.7 & 35.6 & 55.7 & & 0.727$\to$0.785 & 0.759$\to$0.806 \\
 & & published & 4.2 & 47.0 & 48.8 & & & \\
\addlinespace
Riyadh (Durrani 2025)$^\ddagger$ & MSK-2018 & observed & 5.7 & 26.7 & 67.6 & 1496 & & \\
 & & re-weighted & 4.3 & 19.5 & 76.3 (73.8--78.6) & 555 & & \\
 & & LR predicted (weighted) & 6.8 & 34.8 & 58.4 & & 0.715$\to$0.775 & 0.734$\to$0.771 \\
 & & published & 4.2 & 47.0 & 48.8 & & & \\
\addlinespace
Riyadh (Durrani 2025)$^\ddagger$ & Duke MRI & observed & 17.7 & 49.8 & 32.5 & 637 & & \\
 & & re-weighted & 14.0 & 50.1 & 35.9 (32.5--39.7) & 277 & & \\
 & & LR predicted (weighted) & 8.3 & 34.5 & 57.1 & & 0.709$\to$0.710 & 0.779$\to$0.776 \\
 & & published & 4.2 & 47.0 & 48.8 & & & \\
\addlinespace
Pooled Saudi anchors (random effects) & & published & 9.2 & -- & 44.7 (37.2--52.5) & & & \\
\bottomrule
\end{tabular}}
\end{table}

\begin{figure}[htbp]\centering
\adjustbox{max width=\linewidth}{\pgfplotstableread{\gdata{transport}}\tbltr
\begin{tikzpicture}
\begin{axis}[bcbase, xbar stacked, width=0.9\linewidth, height=6.4cm, bar width=8pt, xmin=0, xmax=100,
  xlabel={Patients (\%)}, xmajorgrids, ytick=data, yticklabels from table={\tbltr}{label},
  enlarge y limits=0.07, legend columns=3, legend style={at={(0.5,1.02)}, anchor=south},
  nodes near coords={\pgfmathprintnumber[fixed,precision=0]{\pgfplotspointmeta}},
  nodes near coords style={font=\sffamily\tiny, text=white}, point meta=explicit,
  every node near coord/.append style={anchor=center}]
\addplot[fill=gI, draw=white] table[x=g1, y=y, meta=g1] {\tbltr};
\addplot[fill=gII, draw=white] table[x=g2, y=y, meta=g2] {\tbltr};
\addplot[fill=gIII, draw=white] table[x=g3, y=y, meta=g3] {\tbltr};
\legend{Grade 1, Grade 2, Grade 3}
\draw[inkb, dashed] (axis cs:0,3.5) -- (axis cs:100,3.5);
\end{axis}
\end{tikzpicture}
}
\caption{Grade distributions implied by re-weighting the development and external cohorts to each
Saudi anchor's case-mix, and the published Saudi distributions.}\label{fig:transport}
\end{figure}

\subsection{Sensitivity analysis}
In the Ljubljana cohort (\NLj{} patients, grouped covariates and no receptor data), the best
macro-AUC for ``degree of malignancy'' was \LjBest{} (Supplementary Table~\ref{tab:lj}). This
confirms that discrimination drops when receptor status is not available.
